\section*{THE PARADOXES OF THE THEORY OF SETS AND THE CRISIS OF FOUNDATIONS.}
\phantomsection
\addcontentsline{toc}{section}{THE PARADOXES OF THE THEORY OF SETS AND THE CRISIS OF FOUNDATIONS.}

The first ``paradoxical'' sets appeared in the theory of cardinals and ordinals. In 1897, Burali-Forti remarked that one cannot consider the existence of a set consisting of all the ordinals, because this set would be well ordered, and so isomorphic to one of its segments distinct from the whole, which is absurd {[}42{]}. \(^{55}\) In 1899, Cantor observes (in a letter to Dedekind) that it is no longer possible to speak of the cardinals forming a set, nor to speak of the ``set of all sets'' without reaching a contradiction (the set of subsets of this latter ``set'' \(\Omega\) would be equipotent to a subset of \(\Omega\) , which is contrary to the inequality \(m < 2^{m}\) ) ({[}47{]}, pp.~444-448). In 1905 finally, Russell analysing the proof of this inequality, shows that the reasoning that establishes it proves (without appealing to the theory of cardinals) that the notion of the ``set of all sets which are not elements of themselves'' is, it also, contradictory {[}266{]}. \(^{56}\)

It could be believed that such ``antinomies'' would only be manifested in the peripheral regions of mathematics, characterised by the consideration of sets of a ``size'' inaccessible to intuition. But other ``paradoxes'' were soon to threaten the most classical parts of mathematics. Berry and Russell ({[}266{]}, v. I, pp.63-64), simplifying an argument of J. Richard {[}258{]}, observe in effect that the set of whole numbers of which the definition can be expressed in less than sixteen words in English is finite, but that it is however contradictory to define a whole number as ``the least whole number which is not definable by less than sixteen words in English'', because this definition consists of fifteen words only.

Although such reasoning, so far from that usually used by mathematicians, may have appeared to many of them as a kind of riddle, they pointed the way none the less to the necessity for a revision of the bases of mathematics, aiming to eliminate ``paradoxes'' of this nature. But although there was unanimity concerning the urgency of this revision, radical divergences were soon to appear concerning the manner of realising it. For a first group of mathematicians, be they ``idealists'', be they ``formalists'', \(^{57}\) the situation created by the ``paradoxes'' of the theory of sets was very analogous to that which resulted, in geometry, from the discovery of non-Euclidean geometries or of ``pathological'' curves (such as curves without tangents); it must lead to a similar, but more general, conclusion namely that it is hopeless to seek to base an arbitrary mathematical theory on an appeal (explicit or not) to ``intuition''. The position can be summed up by the very words of the principal adversary of the formalist school: ``The formalist'' says Brouwer ({[}39 a{]}, p.~83), ``submits that human reason does not have at its disposal exact pictures of straight lines or numbers greater than ten, for example \ldots It is true that certain relations between mathematical entities, that we take as axioms, we can deduce from other relations according to fixed rules, with the conviction that in this way we derive truth from other truths by logical reasoning \ldots{[}But{]} for the formalist, mathematical exactitude resides only in the development of the sequence of relations, and is independent of the meaning that one could wish to give to the relations or entities that they link.''

It consists, therefore, for the formalist, in giving to the theory of sets an axiomatic basis quite analogous to that of elementary geometry, where one does not concern oneself with what the ``objects'' are that are called ``sets'', nor what is meant by the relation \(x \in y\) , but where one enumerates the conditions to be imposed on this latter relation; of course, this must be done in such a way as to include, as far as possible, all the results of the theory of Cantor, all the while making it impossible for ``paradoxical'' sets to exist. The first example of such an axiomatisation was given by Zermelo in 1908 {[}342 c{]}; he avoids sets that are ``too large'' by the introduction of an ``axiom of selection (Aussonderung)'', a property \(P(x)\) only determining a set made up of the elements which have it if \(P(x)\) already entails a relation of the form \(x \in A\) .⁵⁸ But the elimination of paradoxes analogous to ``Richard's paradox'' could only be done by restricting the meaning attached to the notion of ``property''; on that, Zermelo contents himself with describing in very vague terms a type of property that he calls ``defined'', and to indicate that one must limit oneself to these latter in the application of the axiom of selection. This point was made precise by Skolem {[}286 a{]} and Fraenkel {[}113 b{]}; as they observed, its elucidation forces us to place ourselves in a completely formalised system, where the notions of ``property'' and of ``relation'' have lost all ``significance'' and have become simple designations for the collections formed following explicit rules. That necessitates of course that one must put into the system the rules of Logic that are used, which was not yet the case in the system of Zermelo-Fraenkel.

Other axiomatisations of the theory of sets were proposed following on from this. Let us mention principally that of von Neumann {[}342 a and c{]} which, more than the system of Zermelo-Fraenkel, approached the primitive conception of Cantor: this latter, to avoid paradoxical sets, had already proposed ({[}47{]}, pp.~445-448), in his correspondence with Dedekind, to distinguish between two kinds of sets, the ``multiplicities'' (``Vielheiten'') and the ``sets'' (``Mengen'') proper, the second ones being characterised in that they can be thought of as a single object. It is this idea that von Neumann makes precise in distinguishing two types of objects, ``sets'' and ``classes''; in his system (almost completely formalised), classes are distinguished from sets in that they cannot be placed on the left of the sign ∈. One of the advantages of such a system is that it rehabilitates the notion of ``universal class'' used by the logicians of the XIXth century (which naturally is not a set); let us point out also that the system of von Neumann avoids (for the theory of sets) the introduction of axiom schemes, replaced by suitable axioms (which makes its logical study easier). Some variants of the system of von Neumann have been given by Bernays and Gödel {[}130 b{]}.

The elimination of paradoxes seems to have been well carried out by the systems above, but at the cost of restrictions which cannot but seem to be very arbitrary. In favour of the system of Zermelo-Fraenkel, one can say that it restricts itself to formulating restrictions that only sanction the current practice in the applications of the notion of set to the various mathematical theories. The systems of von Neumann and of Gödel are further removed from the usual conceptions; on the other hand, it is not excluded that it is easier to insert certain mathematical theories still in their beginnings in the framework provided by such systems rather than in the narrower framework of the system of Zermelo-Fraenkel.

Certainly one could not affirm that any of these solutions give the impression of being definitive. If they satisfy the formalists, it is that these latter refuse to take into consideration the individual psychological reactions of each mathematician; they assume that a formalised language has completed its task when it can transcribe mathematical reasoning into a form lacking ambiguity, and so useful as a vehicle for mathematical thought; everyone is free, they would say, to think what he likes about the ``nature'' of mathematical objects or the ``truth'' of the theorems that he uses, so long as his reasoning can be transcribed into the common language. \(^{59}\)

In other words, from a philosophical point of view, the attitude of the formalists consists in being disinterested in the problem set by the ``paradoxes'', in abandoning the platonic position which aimed to attribute to mathematical notions an intellectual ``content'' common to all mathematicians. Many mathematicians draw back faced with this break with tradition. Russell, for example, seeks to avoid paradoxes by analysing their structure in a deeper way. Taking up again an idea first brought forward by J. Richard (in the article {[}258{]} where he displayed his ``paradox'') and developed later by H. Poincaré {[}251 c{]}, Russell and Whitehead observe that the definitions of paradoxical sets all violate the following principle, called the ``principle of the vicious circle'': ``An element of which the definition involves the totality of the elements of a set can not belong to that set'' ({[}266{]}, v. I, p.~40). Also it is this statement that serves as basis for the Principia, and it is to satisfy it that the ``theory of types'' is developed in this work. Like that of Frege by which it is inspired, the logic of Russell and Whitehead contains ``propositional variables''; the theory of types proceeds to a classification of these different variables, of which the main lines are the following. Starting from a ``domain of individuals'' which are not specified and which can be qualified as ``objects of order 0'', the relationships where the variables (free or tied) are individuals, are said to be ``objects of the first order''; and in general, the relations in which the variables are objects of order ≤ n (one at least being of order n) are said to be ``objects of order \(n + 1\) ''. \(^{60}\) A set of objects of order n can then only be defined by a relation of order \(n + 1\) , a condition that allows the elimination without trouble of the paradoxical sets. \(^{61}\) But the principle of the ``hierarchy of types'' is so restrictive that by adhering to it strictly, one ends up with a mathematics of inextricable complexity. \(^{62}\) To escape this consequence, Russell and Whitehead are constrained to introduce an ``axiom of reducibility'' affirming the existence, for every relation between ``individuals'', of a relation of the first order to which it is equivalent; a condition quite as arbitrary as the axioms of the formalists, and which reduces considerably interest in the construction of the Principia. So the system of Russell and Whitehead has had more success with the logicians than with the mathematicians; it is in any case not entirely formalised, \(^{63}\) and as a result there are numerous obscure details. Various efforts have been made to simplify and clarify this system (Ramsey, Chwistek, Quine, Rosser); tending to use languages that are more and more fully formalised, these authors replace the rules of the Principia (which still had a certain intuitive foundation) with restrictions only taking account of the wording of the aggregations being considered; not only do these rules appear to be as gratuitous as the prohibitions formulated in the systems of Zermelo-Fraenkel or of von Neumann, but, being further removed from the practice of mathematicians, they have, in many cases, led to unacceptable consequences, which the author had not foreseen (such as the paradox of Burali-Forti, or the negation of the axiom of choice).

For mathematicians of the preceding schools, it was above all a case of not renouncing any part of the heritage of the past: ``From the paradise that Cantor has created for us'', says Hilbert ({[}163 c{]}, p.~274), ``no one should be able to chase us''. To do this, they are willing to accept limitations on mathematical reasoning, hardly essential because consonant with use, but which do not seem to be imposed by our mental customs and the intuitive idea of the notion of set. Anything seems to them preferable to the intrusion of psychology into the criteria of validity in mathematics; and rather than ``putting into account the characteristics of our brains'', as Hadamard says ({[}12{]}, p.~270), they are resigned to imposing on the mathematical domain limits which are for the most part arbitrary, so long as they include classical Mathematics and do not risk hampering future progress.

It is a very different matter when we see the attitude of mathematicians belonging to the tendency of which it remains for us to speak. If the formalists accept the renunciation of the control of the ``eyes of the mind'' so far as mathematical reasoning is concerned, the mathematicians that have been called ``empiricists'', ``realists'' or ``intuitionists'' refuse this abdication; they must have a kind of internal certainty guaranteeing the ``existence'' of the mathematical objects with which they are concerned. As long as it was only a case of renouncing spatial intuition, there was no serious objection, since the arithmetical ``models'' allowed retrenchment behind the intuitive notion of whole numbers. But irreducible opposition manifests itself when it is a matter of reducing the notion of whole number to that (much less intuitively precise) of set, then of imposing on the handling of sets barriers without intuitive foundation. The first in time of these opponents (and the one who, by the authority of his genius, was to exercise the most influence) was H. Poincaré; having admitted, not only the axiomatic point of view for geometry and the arithmetisation of Analysis, but also a good part of the theory of Cantor (which he was one of the first to apply fruitfully in his work), he refuses on the other hand to conceive that arithmetic may be able, it also, to be justified by an axiomatic treatment; the principle of complete induction seems to him in particular a fundamental intuition of our mind, where it is impossible to see a pure convention \(^{64}\) {[}251 c{]}. Hostile in principle to formalised languages, of which he contested the usefulness, he confuses constantly the notion of whole number in formalised mathematics, and the use of whole numbers in the theory of proof, which was starting up then, and of which we will speak later; no doubt it was difficult at that time to make as clearly as today --- after 50 years of study and discussion --- this distinction which however had been clearly felt by a Hilbert or a Russell.

The critics of this kind multiplied after the introduction of the axiom of choice by Zermelo in 1904 {[}342 a{]}. Its use in many previous proofs in Analysis or the theory of sets had until then gone by almost unnoticed; \(^{65}\) it is in following an idea suggested by Erhard Schmidt that Zermelo, explicitly formulating this axiom, deduced from it in an ingenious way a satisfactory proof of the existence of a well ordering on every set. Coming at the same time as the ``paradoxes'', it appears that this new way of reasoning, by its odd allure, had thrown many a mathematician into confusion; it is enough to see the strange misunderstandings which surged up about this topic, in the following volume of Mathematische Annalen, under the signature of mathematicians who were as familiar with the Cantorian methods as Schoenflies and F. Bernstein. The criticisms of E. Borel, published in this same volume, are more substantial and attach themselves clearly to the point of view expressed by H. Poincaré about whole numbers; they are developed and discussed in an exchange of letters between E. Borel, Baire, Hadamard and Lebesgue, who had remained in the classical tradition in French mathematics {[}12{]}. Borel starts by denying the validity of the axiom of choice because it contains in general an uncountable infinity of choices, which is inconceivable intuitively. But Hadamard and Lebesgue observe that a countable infinity of arbitrary successive choices is no more intuitive, since it consists of an infinity of operations, that it is impossible to conceive as really taking place. For Lebesgue, who enlarged the debate, everything reduces to knowing what is meant by saying that a mathematical object ``exists''; it is necessary, for him, that one ``names'' explicitly a property defining it in a unique way (a ``functional'' property, we would say); for a function such as the one used by Zermelo in his reasoning, that is what Lebesgue calls a ``law'' of choice; if, he continues, this requirement is not satisfied, and that ones limits oneself to ``thinking'' about this function instead of ``naming'' it, can one be sure that during the reasoning process one is always thinking of the same thing ({[}12{]}, p.~267)? In any case, this brings Lebesgue to new doubts; already the question of the choice of a single element in a set seems to him to raise difficulties: one must be sure that such an element ``exists'', that is to say that one can ``name'' at least one of the elements of the set. \(^{66}\) Can one then speak of the ``existence'' of a set of which one does not know how to ``name'' each element? Already Baire does not hesitate to deny the ``existence'' of the set of all subsets of a given infinite set (loc. cit., pp.~263-264); in vain does Hadamard observe that these requirements lead to renouncing even the possibility of speaking of the real numbers: it is to just that conclusion that E. Borel finally agrees. Putting aside the fact that countability seems to have acquired freehold, one has returned just about to the classical position of the adversaries of the ``actual infinite''.

All these objections were not very systematic; it was down to Brouwer and his school to undertake a complete remould of mathematics guided by similar, but even more radical, principles. We would not dare to summarise here a doctrine as complex as intuitionism, which consists as much of psychology as mathematics, and we will limit ourselves to indicating some of the most striking aspects, referring for more details to the work of Brouwer himself {[}39 a and b{]} and to the exposition by Heyting {[}162{]}. For Brouwer, Mathematics is identical with the ``exact'' part of our thought, based on the first intuition of the sequence of natural numbers, and which it is impossible to translate without mutilation into a formal system. It is in any case only ``exact'' in the thought of mathematicians, and it is a chimera to hope to develop an instrument of communication between them which is not subject to all the imperfections and ambiguities of language; at most it is possible to hope to awaken in the interlocutor a favourable state of mind by descriptions that are more or less vague ({[}162{]}, pp.~11-13). Intuitionist mathematics attaches hardly more importance to logic than to language: a proof is not conclusive by virtue of logical rules fixed once and for all, but because of the ``immediate evidence'' of each of its steps. This ``evidence'' must moreover be interpreted in a more restrictive way than by E. Borel and his followers: it is thus that in intuitionist mathematics, it cannot be said that a relation of the form ``R or (not R)'' is true (principle of the excluded middle), unless for all system of values given to the variables appearing in R, one can prove that one of the two propositions R, ``not R'' is true; for example, from the equation ab = 0 between two real numbers, it cannot be concluded that ``a = 0 or b = 0'', for it is easy to create explicit examples of real numbers a, b for which one has ab = 0 without knowing, at the actual time, how to prove either of the two propositions a = 0, b = 0 ({[}162{]}, p.~21).

It is not astonishing that, starting from such principles, the intuitionist mathematicians finish up with results that are very different from the classical theorems. A whole part of these disappear, for example most of the ``existential'' theorems in Analysis (such as the theorems of Bolzano and Weierstrass for real functions); if a function of a real variable ``exists'' in the intuitionist sense, it is ipso facto continuous; a bounded monotone sequence of real numbers does not necessarily have a limit. On the other hand, many classical notions become ramified, for the intuitionist, into several distinct fundamental notions: there are thus two notions of convergence (for a sequence of real numbers) and eight of countability. It goes without saying that transfinite induction and its applications to modern Analysis are (as for the greater part of the theory of Cantor) condemned without appeal.

It is only in this way that, according to Brouwer, mathematical propositions can acquire a ``content''; the formalist reasonings that go beyond what the intuitionists allow are judged to be without value, since one can no longer give them a ``meaning'' to which the intuitive notion of ``truth'' could be applied. It is clear that such judgements can only rely on a pre-existing notion of ``truth'', of a psychological or metaphysical nature. It is to say in practice that they escape all discussion.

It is undoubted that the vigorous attacks that came from the intuitionist camp have sometimes forced not only the mathematical schools in the avant-garde, but even the followers of traditional mathematics to go on the defensive. A famous mathematician has admitted having been impressed by the attacks to the point of having voluntarily restricted his work to branches of mathematics judged to be ``certain''. But such cases can hardly have been frequent. The intuitionist school, of which the memory is no doubt destined to remain only as an historical curiosity, would at least have been of service by having forced its adversaries, that is to say definitely the immense majority of mathematicians, to make their position precise and to take more clearly notice of the reasons (the ones of a logical kind, the others of a sentimental kind) for their confidence in mathematics.

